%%
%% Director-CenterNet: Multi-Region Viewport Prediction for Automated Esports Broadcasting
%% Target Journal: Expert Systems with Applications (Elsevier)
%% Template: elsarticle-template-num.tex (numbered citation style)
%%

\documentclass[preprint,12pt]{elsarticle}

%% Use the option review to obtain double line spacing
%% \documentclass[authoryear,preprint,review,12pt]{elsarticle}

%% Use the options 1p,twocolumn; 3p; 3p,twocolumn; 5p; or 5p,twocolumn
%% for a journal layout:
%% \documentclass[final,1p,times]{elsarticle}
%% \documentclass[final,3p,times]{elsarticle}
%% \documentclass[final,5p,times]{elsarticle}

%% The amssymb package provides various useful mathematical symbols
\usepackage{amssymb}
%% The amsmath package provides various useful equation environments.
\usepackage{amsmath}
%% Extended theorem environments
%% \usepackage{amsthm}

%% Line numbers for review
%% \usepackage{lineno}

%% Additional packages
\usepackage{booktabs}
\usepackage{multirow}
\usepackage{algorithm}
\usepackage{algorithmic}
\usepackage{xcolor}
\usepackage{subcaption}
\usepackage{mathtools}

%% Convenience macros
\newcommand{\R}{\mathbb{R}}
\newcommand{\todo}[1]{\textcolor{red}{[\textbf{TODO}: #1]}}

\journal{Expert Systems with Applications}

\begin{document}

\begin{frontmatter}

%% Title
\title{Director-CenterNet: Multi-Region Viewport Prediction\\
       for Automated Esports Broadcasting}

%% Authors
%% use the corref command within \author for corresponding author footnote
\author[gist]{Cheong-mok Bae}
\ead{cmbae0307@gm.gist.ac.kr}

\author[gist,aigist]{Kyung-Joong Kim\corref{cor1}}
\ead{kjkim@gist.ac.kr}

\cortext[cor1]{Corresponding author}

\affiliation[gist]{organization={School of Integrated Technology},
             addressline={123 Cheomdan-gwagi-ro, Buk-gu},
             city={Gwangju},
             postcode={61005},
             country={South Korea}}

\affiliation[aigist]{organization={AI Graduate School},
             addressline={123 Cheomdan-gwagi-ro, Buk-gu},
             city={Gwangju},
             postcode={61005},
             country={South Korea}}

%% Abstract
\begin{abstract}
Human game observers, who control in-game cameras and provide viewers with an
engaging experience, are a vital part of electronic sports (Esports). In
professional broadcasts multiple observers cooperate to operate several camera
feeds at once, because a single view cannot cover the events that occur
concurrently across a large map. Consequently, various methods to create
automatic observers have been explored. However, these methods predict only a
single region per frame. This is problematic because the attention of human
observers is not unimodal: observers watching the same replay agree on only
about half of their viewport area, and their attention separates further as the
game progresses. Previous learning-based observers treated this disagreement as
noise and collapsed it into a single consensus region, which discards the
minority modes and causes the predicted camera to oscillate between competing
regions across consecutive frames. In this paper, we overcome these problems by
reformulating automatic observing as \emph{Multi-Region Viewport Prediction}
(MRVP) and solving it with a heatmap-based detection method, CenterNet, in a
real-time strategy game (e.g., StarCraft). Instead of collapsing the modes of
human attention, the proposed model, \emph{Director-CenterNet}, recovers them:
we rank the modes of the aggregated observer distribution by the number of
observers they represent, and assign the dominant mode to a primary region and
the remaining modes to auxiliary regions. The predicted regions are
presentation-agnostic, so they can be rendered as a main screen with
picture-in-picture inserts, as a split-screen layout, or as a ranked candidate
set for a human director. We design four training objectives for this task: a
\emph{Human Consensus Match} loss, a \emph{Ranked Mode Coverage} loss, a
\emph{Spatial Repulsion} loss, and a \emph{Trajectory Smoothness} loss.
We further show that a detector-based observer already emits several candidate
regions per frame, so the contribution of the proposed formulation is not
multiplicity as such but that the regions are ranked by observer support,
non-redundant, and temporally stable — which is what determines how many of the
disagreeing observers the output actually serves. Controlling a proposal
detector's region count by its confidence threshold, we find that the number of
regions it emits varies inversely with the number of modes the frame contains,
at every threshold and with no sign change available, because a per-proposal
confidence is spread thinner precisely when the evidence is divided. The
heatmap formulation reverses this dependence. At a matched region budget the
proposed method leads on observer coverage and on count calibration; given an
unconstrained budget the strongest single-region baseline remains ahead on
primary-region accuracy, and we report both.
\end{abstract}

%% Graphical abstract
\begin{graphicalabstract}
%% \includegraphics{graphical_abstract}
\todo{Insert graphical abstract: pipeline diagram showing game-state tensor
$\rightarrow$ Director-CenterNet $\rightarrow$ Main + PIP viewports overlaid
on minimap.}
\end{graphicalabstract}

%% Research highlights
\begin{highlights}
\item Human observer attention in real-time strategy games is shown to be
      multimodal, not unimodal.
\item Automatic observing is reformulated as multi-region viewport prediction
      for Esports broadcasting.
\item Minority observer attention, previously discarded as noise, becomes
      supervision for auxiliary regions.
\item Repulsion and smoothness objectives remove redundant overlap and camera
      thrashing.
\item Predicted regions are presentation-agnostic and support PIP,
      split-screen, and director assistance.
\end{highlights}

%% Keywords
\begin{keyword}
Esports \sep Automatic observer \sep Multi-viewport prediction \sep
CenterNet \sep Saliency-guided learning \sep Spatial resource allocation \sep
StarCraft \sep Behavior cloning
\end{keyword}

\end{frontmatter}

%% \linenumbers   %% uncomment for review line numbers

%%====================================================================
\section{Introduction}
\label{sec:intro}
%%====================================================================

Human game observers in electronic sports (Esports) fulfill a dual role: 
they must simultaneously identify the most strategically significant region 
of the map \emph{and} coordinate multiple camera feeds to ensure that 
concurrent events are not missed \cite{joo2023learning}. In professional 
Esports productions, a team of observers typically manages one main broadcast 
camera together with several auxiliary picture-in-picture (PIP) channels, each 
covering a distinct map region. Automating this workflow is challenging for 
three reasons: (1) the number of concurrent high-priority events is dynamic; 
(2) the spatial relationships among events evolve rapidly; and (3) the 
subjective notion of ``broadcast saliency'' cannot be reduced to a single, 
well-defined objective \cite{bae2025observer}.

Prior work on automatic Esports observers has focused on the 
\emph{single-viewport} setting, where the goal is to select one rectangular 
window of fixed size from the game map at each frame. Rule-based approaches 
\cite{mattsson2015automatic} define priority scores over hand-crafted events 
and switch the camera accordingly. Learning-based methods 
\cite{joo2023learning} instead collect observational data from human
spectators and train a detector-style model with Region of Common Interest
(ROCI) loss to imitate consensus human camera behaviour. Subsequent work added
an auxiliary objective derived from fixed-kernel spatial priors, which improved
agreement with aggregated human attention without changing inference-time
outputs \cite{anonymous2025kbrs}.

Despite these advances, we argue that the single-region setting is not merely
restrictive but \emph{mismatched to the supervision it is trained on}. The
evidence for this mismatch is already present in previous studies. Observers
watching the same replay share only about 52\% of their viewport area with each
other, and questionnaires show that they attend to different aspects of the
match, such as building upgrades, resource expansion, or fighting
\cite{joo2023learning}. Moreover, this disagreement is not constant over a
match: it varies with what is happening on the map, and so does the accuracy of
every method built on it. We therefore regard the roughly 50\% agreement
ceiling not as measurement noise to be averaged away, but as evidence that the
target distribution is \emph{multimodal}.

Section~\ref{subsec:target_structure} quantifies this directly. Extracting the
modes of the aggregated observer distribution and measuring a strong
single-region observer against them, we find that $59\%$ of frames contain two
or more modes, and that the observer's accuracy is governed by how many. IR
falls monotonically from $0.779$ on unimodal frames to $0.499$ on frames with
five modes, while the fraction of frames whose predicted region lands on a
\emph{minority} mode rather than the consensus one rises from $0$ to $0.48$.
The same relation holds when the modes are ranked by support: when the top two
modes are supported by equal numbers of observers, IR is $0.595$ and the
prediction alternates between them on $7.1\%$ of consecutive frame pairs; when
a single mode is unanimous, IR is $0.853$ and that alternation rate is exactly
zero.

We stress that the controlling variable is the number of modes, not elapsed
time. Accuracy on our test set is not monotonic in game progression: it is
highest early, lowest at the midpoint, and highest of all at the end, tracking
a mode count that peaks in mid-game ($2.60$) and falls as the match resolves
into a single decisive engagement ($1.87$). Game progression is a proxy for
multimodality, and it is an imperfect one.

This interpretation explains the limitations of previous work. ROCI
\cite{joo2023learning} and fixed-kernel auxiliary priors
\cite{anonymous2025kbrs} are both devices for rescuing a single prediction on a
multimodal target: they sharpen or re-weight the distribution so that one
region can be selected from it. Such devices necessarily discard the minority
modes, and they become least effective in exactly the phase in
which the modes separate. In addition, when two regions receive nearly equal
support, a model trained to select one of them oscillates between them across
consecutive frames.

We address both limitations by reformulating automatic observing as
\textbf{Multi-Region Viewport Prediction (MRVP)}: at each frame, the observer
produces an ordered set of $K$ non-redundant regions of fixed size. Unlike the
single-region approaches mentioned previously, this formulation does not
collapse the modes of human attention but covers them, and it allows camera
thrashing to be suppressed explicitly rather than implicitly. We emphasise that
the predicted regions are \emph{presentation-agnostic}. The same output can be
rendered as a main screen with picture-in-picture inserts, as a split-screen or
multi-view grid, or as a ranked candidate set presented to a human director in
the manner of \cite{chen2018learning}. The prediction problem is therefore
independent of the production layout eventually chosen.

Furthermore, we propose to obtain the supervision for the auxiliary regions
from the human observational data itself, rather than from a hand-designed
prior. Because the aggregated observer distribution already contains the
secondary modes that previous work suppressed, we smooth this distribution,
extract its local maxima, and rank the resulting modes by the number of
observers they represent. The dominant mode supervises the primary region and
the remaining modes supervise the auxiliary regions. In other words, the
minority attention that ROCI treated as an obstacle becomes the training signal
for the additional regions.

To solve MRVP we propose \textbf{Director-CenterNet}, built on CenterNet 
\cite{zhou2019objects}. CenterNet's dense heatmap representation is 
well-suited to MRVP because (i) it simultaneously scores all candidate 
viewport centres, enabling multi-peak extraction without iterative proposal 
generation; (ii) the dense heatmap preserves the multimodal structure of the
target, whereas an instance mask requires a discrete assignment that discards
it; and (iii) local-maximum extraction naturally supports a variable number of
regions, allowing an auxiliary region to be suppressed when no secondary mode
is detected.

Our contributions are as follows.

\begin{enumerate}
  \item We show that human observer attention in a real-time strategy game is
        multimodal and that its modes separate as the match progresses, and we
        argue that single-region observers are therefore mismatched to their
        own supervision (Section~\ref{sec:problem}).
  \item We formalise Multi-Region Viewport Prediction (MRVP) as a
        presentation-agnostic task and propose evaluation metrics for
        multi-region coverage, per-observer service, and temporal stability
        (Section~\ref{sec:problem}).
  \item We propose Director-CenterNet, a CenterNet-based multi-peak observer
        that predicts an ordered set of regions in a single forward pass,
        together with a Ranked Mode Coverage loss that supervises the auxiliary
        regions using minority observer attention rather than a hand-designed
        prior (Section~\ref{subsec:rmc_loss}).
  \item We show that the baselines already emit several candidate regions per
        frame, and evaluate them on their own top-$K$ accordingly, so that the
        comparison isolates whether the regions are ranked, non-redundant and
        stable rather than merely numerous. Controlling their region count by
        the only mechanism they offer — the confidence threshold — we find that
        the number of regions a proposal detector emits varies
        \emph{inversely} with the number of modes the frame actually contains,
        at every threshold and with no sign change available, whereas the
        heatmap formulation reverses that dependence. At a matched region
        budget the proposed method is ahead on coverage and on count
        calibration; given an unconstrained budget the Mask R-CNN baseline
        remains ahead on single-region accuracy, and we report both
        (Section~\ref{sec:experiments}).
\end{enumerate}

%%====================================================================
\section{Background and Related Work}
\label{sec:related}
%%====================================================================

\subsection{Single-Viewport Automatic Observers}

Automatic game observation has been studied primarily in real-time strategy 
(RTS) games, where the viewport selection problem is most challenging due to 
large map size and high event concurrency \cite{ontanon2013survey}. 
Rule-based observers such as AIIDE and SSCAIT \cite{mattsson2015automatic} 
define priority scores over game events and teleport the camera to the 
highest-priority event at each decision step. While interpretable, these 
systems require manual event taxonomies and are brittle under novel strategies.

Joo et al. \cite{joo2023learning} proposed the first event-free learning-based 
observer for StarCraft. They framed viewport selection as an instance 
segmentation problem, treating the human observer's viewport as the object to 
detect. Using Mask R-CNN \cite{he2020mask} and a Region of Common Interest 
(ROCI) loss that amplifies the overlapping area among multiple human observers, 
their method outperforms both rule-based baselines and behaviour cloning. 
Subsequent work added an auxiliary objective derived from fixed-kernel spatial
priors, which improved agreement with aggregated human attention without
changing inference-time outputs \cite{anonymous2025kbrs}. Both studies report
that accuracy declines as the match progresses and that the viewports of
observers become divided in the later phases — the observation this work builds
on.

The Observer framework \cite{bae2025observer} provides an open-source, 
reproducible pipeline for data collection, preprocessing, and model training 
on StarCraft: Brood War replays, which we use in this work.

\subsection{Object Detection Architectures}

Our work builds on two complementary detection paradigms. 
\textbf{CenterNet} \cite{zhou2019objects} frames detection as keypoint 
heatmap regression: the network predicts a Gaussian heatmap peaked at each 
object's centre, with offset and size regression heads for bounding box 
recovery. Multi-peak extraction by local-maximum search provides a natural 
mechanism for detecting a variable number of objects per image. 
\textbf{DETR} \cite{carion2020end} and RT-DETR \cite{lv2023detrs} reframe 
detection as set prediction via Hungarian matching, eliminating 
non-maximum suppression. We choose CenterNet over transformer-based detectors 
for MRVP because the heatmap formulation preserves the multimodal structure of
our target and admits a variable number of outputs without modifying a query
set.

\subsection{Saliency and Multi-Region Detection}

The problem of distributing a limited number of viewports over a dynamic scene 
relates to multi-salient region detection \cite{shi2016hierarchical}, active 
perception \cite{chen2011active}, and sensor placement 
\cite{krause2008near}. Classical saliency models incorporate hand-designed 
priors — including centre bias and background priors — to guide attentional 
selection \cite{itti1998model, wei2012geodesic, tatler2007central}. Our work 
differs in that salient regions are defined by aggregated human attention over 
a structured multi-channel game-state tensor rather than by low-level visual 
features in natural images.

\subsection{Temporal Stability in Video-Based Observers}

Camera thrashing — rapid oscillation between competing salient regions — is a 
known failure mode in learning-based observers. Trajectory smoothness losses 
penalising inter-frame displacement have been applied in video object detection 
\cite{feichtenhofer2017detect} and automated cinematography 
\cite{gandhi2014video} to reduce temporal inconsistency. We adapt this 
approach to the MRVP setting by penalising the L2 displacement of predicted 
viewport centres across consecutive frames.

%%====================================================================
\section{Problem Formulation}
\label{sec:problem}
%%====================================================================

\subsection{Input Representation}

Following \cite{joo2023learning, bae2025observer, anonymous2025kbrs}, at 
each game frame $t$ we construct a multi-channel state tensor 
$\mathbf{X}_t \in \mathbb{R}^{C \times H \times W}$ on a tile grid of size 
$H \times W$. Channels encode semantically distinct entity categories for 
each player (workers, ground units, air units, buildings), terrain elevation, 
combined vision, and resource locations. Each tile corresponds to 
$32 \times 32$ pixels in the game world.

To incorporate short-term temporal context, we stack frames at a fixed 
stride $\Delta$:
\begin{equation}
  \mathbf{X}^{\text{stk}}_t 
    = \text{concat}\bigl[\mathbf{X}_t,\,\mathbf{X}_{t-\Delta},\,
      \ldots,\,\mathbf{X}_{t-W\Delta}\bigr]
    \in \mathbb{R}^{(W+1)C \times H \times W}.
  \label{eq:stacking}
\end{equation}
Unless otherwise stated we use $\Delta = 8$ and $W = 3$, yielding a 
four-frame stack.

\subsection{Human Supervision and Aggregation}

For each frame $t$ and human observer $u \in \{1,\ldots,U\}$, we record a 
viewport $V^{(u)}_t$ specified by its top-left tile coordinate 
$\mathbf{p}^{(u)}_t \in \mathcal{C}$, where
\begin{equation}
  \mathcal{C} = \{(i,j) \mid 0 \le i \le H - h,\; 0 \le j \le W - w\}
  \label{eq:feasible}
\end{equation}
is the set of feasible top-left positions for a viewport of fixed tile size 
$(h, w)$. We aggregate human viewports into a hard union set
\begin{equation}
  \mathcal{H}_t = \bigcup_{u=1}^{U} V^{(u)}_t
  \label{eq:union}
\end{equation}
and a soft crowd coverage map
\begin{equation}
  a_t(\mathbf{q}) = \frac{1}{U} \sum_{u=1}^{U} 
    \mathbf{1}\!\bigl[\mathbf{q} \in V^{(u)}_t\bigr],
  \label{eq:coverage}
\end{equation}
which measures the fraction of observers whose viewport covers tile
$\mathbf{q}$.

\subsection{Ranked Modes of the Observer Distribution}
\label{subsec:modes}

Rather than reducing $a_t$ to a single consensus point, we recover its modes.
We smooth the coverage map with a Gaussian filter, $\tilde{a}_t = G_{\sigma_m}
* a_t$, and take its local maxima above a relative floor
$\theta \cdot \max \tilde{a}_t$. Candidates are visited in descending order of
$\tilde{a}_t$ and greedily suppressed within a minimum separation $D$, which
prevents a single broad region of agreement from producing several modes. Each
surviving mode $\mathbf{m}$ is assigned a \emph{support}, the number of
observers whose viewport covers it:
\begin{equation}
  n_t(\mathbf{m}) = \sum_{u=1}^{U} \mathbf{1}\bigl[\mathbf{m} \in V^{(u)}_t\bigr].
  \label{eq:support}
\end{equation}
Sorting by support in descending order gives the ordered sequence
\begin{equation}
  \mathcal{M}_t = \bigl(\mathbf{m}^{(1)}_t, \mathbf{m}^{(2)}_t, \ldots\bigr),
  \qquad n_t(\mathbf{m}^{(1)}_t) \ge n_t(\mathbf{m}^{(2)}_t) \ge \cdots,
  \label{eq:ranked_modes}
\end{equation}
which we refer to as the Top-1, Top-2, \ldots\ regions of human attention at
frame $t$. The Top-1 region corresponds to the consensus region targeted by
previous single-region work; the remaining entries are the minority modes that
previous work suppressed.

Two properties of $\mathcal{M}_t$ matter for this study. First, its cardinality
varies over a match, and Section~\ref{subsec:target_structure} shows that it,
rather than elapsed time, is what governs single-region accuracy. Second, the
resolution of the ranking is bounded by $U$. We inherit the observational data
of \cite{joo2023learning} unchanged, which records $U=5$ observers per replay,
so a mode's support is an integer in $\{1,\ldots,5\}$ and by
Eq.~\ref{eq:gt_heatmap} the smallest representable amplitude is $1/U = 0.2$.
A mode carried by a single observer is therefore at the edge of what five
samples can distinguish from idiosyncratic attention. The evidence that the
target is genuinely multimodal does not rest on those single-observer modes:
$18\%$ of frames have their top two modes supported by \emph{equal} numbers of
observers (Table~\ref{tab:by_margin}), which a larger pool would refine but
which is not an artefact of a small one. We return to this bound in
Section~\ref{sec:discussion}.

\subsection{Multi-Region Viewport Prediction (MRVP)}

\textbf{Definition.} Given a game-state tensor $\mathbf{X}^{\text{stk}}_t$, 
predict a set of $K$ viewport positions 
$\{\hat{\mathbf{p}}^{(k)}_t\}_{k=1}^{K} \subset \mathcal{C}$ such that:
\begin{enumerate}
  \item $\hat{\mathbf{p}}^{(1)}_t$ (main viewport) maximally overlaps the 
        human consensus region $\mathcal{H}_t$;
  \item $\hat{\mathbf{p}}^{(k)}_t$ for $k \ge 2$ (PIP viewports) cover 
        salient regions not covered by the main viewport;
  \item Predicted viewports are spatially diverse (non-redundant);
  \item Predicted viewport positions vary smoothly across consecutive frames.
\end{enumerate}

This formulation generalises single-viewport prediction (the special case 
$K=1$) while connecting MRVP to the multi-salient region detection literature.

\subsection{Evaluation Metrics}

\paragraph{Intersection Ratio (IR).}
\begin{equation}
  \text{IR}_t = \frac{\bigl|V(\hat{\mathbf{p}}^{(1)}_t) \cap \mathcal{H}_t\bigr|}
                     {\bigl|V(\hat{\mathbf{p}}^{(1)}_t)\bigr|}
  \label{eq:ir}
\end{equation}

\paragraph{Intersection@$\delta$.}
\begin{equation}
  \text{I@}\delta_t = \mathbf{1}[\text{IR}_t \ge \delta],\quad
  \delta \in \{0,\, 0.3,\, 0.5\}.
  \label{eq:iat}
\end{equation}

\paragraph{Best-of-$K$ Intersection@$\delta$ (BoK@$\delta$).}
Whether \emph{any} predicted region reaches overlap $\delta$, which is the
multi-region analogue of the first-versus-second-suggestion comparison reported
in \cite{chen2018learning}:
\begin{equation}
  \text{BoK@}\delta_t = \mathbf{1}\!\left[\max_{k \le K}
    \frac{\bigl|V(\hat{\mathbf{p}}^{(k)}_t) \cap \mathcal{H}_t\bigr|}
         {\bigl|V(\hat{\mathbf{p}}^{(k)}_t)\bigr|} \ge \delta\right].
  \label{eq:bok}
\end{equation}

\paragraph{Observer Coverage (OC@$\delta$).}
The fraction of \emph{individual} observers served by at least one predicted
region, which measures how many of the disagreeing observers the output
satisfies:
\begin{equation}
  \text{OC@}\delta_t = \frac{1}{U}\sum_{u=1}^{U} \mathbf{1}\!\left[
    \max_{k \le K} \frac{\bigl|V(\hat{\mathbf{p}}^{(k)}_t) \cap V^{(u)}_t\bigr|}
                        {\bigl|V^{(u)}_t\bigr|} \ge \delta\right].
  \label{eq:oc}
\end{equation}
Note that $\text{IR}$ and $\text{I@}\delta$ are computed on the primary region
only and are therefore directly comparable to previously reported
single-region results, whereas BoK@$\delta$ and OC@$\delta$ characterise the
multi-region setting. The gap between $\text{I@}\delta$ and BoK@$\delta$
indicates how much of the human attention distribution is reachable by
additional regions; OC@$\delta$ indicates how many of the disagreeing observers
the output actually serves.

A remark on how the baselines are scored. Detector-based single-region
observers do not emit one box per frame; they emit a ranked list and previous
work reports only its top entry. We therefore evaluate BoK@$\delta$ and
OC@$\delta$ on their top-$K$ boxes as well, rather than treating them as
producing a single region by construction. This is a deliberately demanding
comparison: it credits the baselines with multi-region behaviour they were
never trained for, and it raises the bar substantially — Mask R-CNN + prior
reaches OC3@0.5 $=0.592$ against its own OC1@0.5 of $0.371$. The distinction
the proposed method must establish is consequently not that it produces several
regions, since a detector already does, but that the regions it produces are
ranked by observer support, non-redundant, and temporally stable.

\paragraph{Count calibration.}
Let $\hat{n}_t = |\hat{\mathcal{P}}_t|$ be the number of regions the model
emits at frame $t$ and $n_t = |\mathcal{M}_t|$ the number of ranked modes the
observers produced. Three quantities summarise how the former tracks the
latter. The first two are its accuracy,
\begin{equation}
  \text{MAE}_n = \frac{1}{T}\sum_t \bigl|\hat{n}_t - n_t\bigr|,
  \qquad
  \text{Exact}_n = \frac{1}{T}\sum_t \mathbf{1}\bigl[\hat{n}_t = n_t\bigr],
  \label{eq:count_err}
\end{equation}
and the third is its \emph{responsiveness}: the slope $\beta_n$ of the weighted
least-squares fit of the conditional mean $\mathbb{E}[\hat{n}_t \mid n_t = n]$
on $n$,
\begin{equation}
  \beta_n = \frac{\sum_{n} w_n \bigl(n - \bar{n}\bigr)
                  \bigl(\mathbb{E}[\hat{n}_t \mid n_t = n] - \bar{\hat{n}}\bigr)}
                 {\sum_{n} w_n \bigl(n - \bar{n}\bigr)^2},
  \qquad w_n = \bigl|\{t : n_t = n\}\bigr|.
  \label{eq:count_slope}
\end{equation}
The two are independent. $\text{MAE}_n$ is a level: any model can be brought
close to the mean of $n_t$ by tuning how freely it emits regions, and a
constant predictor $\hat{n}_t \equiv \mathbb{E}[n_t]$ minimises it among all
constant predictors. $\beta_n$ is a direction: it is zero for every constant
predictor regardless of level, positive only when the model emits more regions
on frames whose target genuinely has more modes, and negative when it does the
opposite. Section~\ref{subsec:budget} shows that these two come apart sharply
in practice, and that only the second distinguishes the formulations.

\paragraph{Frame coverage.}
The fraction of evaluated frames for which the model emits at least one region
above its confidence threshold. All metrics above are conditional on the model
having produced an output, so a method that declines to answer on part of the
test set is not comparable to one that answers everywhere unless coverage is
reported alongside. Where a comparison spans methods with different coverage we
report the coverage-penalised value, which scores an unanswered frame as $0$,
or restrict all methods to the frames they all answer; both are stated
explicitly where used.

\paragraph{Viewport Displacement (VD).}
Average L2 displacement of the main viewport centre across consecutive
frames, measuring temporal stability:
\begin{equation}
  \text{VD} = \frac{1}{T-1}\sum_{t=1}^{T-1}
    \bigl\|\bar{\mathbf{p}}^{(1)}_t - \bar{\mathbf{p}}^{(1)}_{t+1}\bigr\|_2,
  \label{eq:vd}
\end{equation}
where $\bar{\mathbf{p}}^{(k)}_t$ denotes the centre of the $k$-th predicted 
viewport at frame $t$.

%%====================================================================
\section{Director-CenterNet}
\label{sec:method}
%%====================================================================

\subsection{Architecture Overview}

Director-CenterNet is a heatmap-based multi-peak detector that extends 
CenterNet \cite{zhou2019objects} to the MRVP problem. The architecture 
comprises three components: (1) a ResNet-FPN backbone encoder that transforms 
the stacked game-state tensor into a spatial feature map; (2) CenterNet heads 
that regress a dense Gaussian heatmap of viewport centres together with offset 
and size maps; and (3) a ranked multi-peak extraction step. 
Fig.~\ref{fig:architecture} provides an overview.

\begin{figure}[t]
  \centering
  \todo{Insert architecture figure: stacked game-state tensor $\rightarrow$
  ResNet-FPN backbone $\rightarrow$ heatmap / offset / size heads
  $\rightarrow$ ranked multi-peak extraction $\rightarrow$ ordered region set,
  with the four training objectives annotated on the ranked human attention
  modes.}
  \caption{Overview of Director-CenterNet. The stacked game-state tensor is
  processed by a ResNet-FPN backbone. The heatmap head produces a dense score
  map over all candidate region centres, and ranked multi-peak extraction
  yields an ordered set of regions in a single forward pass. Training is
  supervised by the ranked modes of the aggregated observer distribution
  (Section~\ref{subsec:modes}): the Top-1 mode supervises the primary region
  and the remaining modes supervise the auxiliary regions.}
  \label{fig:architecture}
\end{figure}

\subsection{CenterNet Backbone and Heads}

\paragraph{Input.} The stacked tensor
$\mathbf{X}^{\text{stk}}_t \in \mathbb{R}^{(W+1)C \times H \times W}$ is
fed to a ResNet-50 \cite{he2016deep} backbone with a Feature Pyramid Network
(FPN) \cite{lin2017feature} to obtain a feature map
$\mathbf{F}_t \in \mathbb{R}^{C_f \times H_s \times W_s}$ at stride $s$.
The tile map is consumed at its native resolution without resizing, and we
take the finest FPN level, $s=4$. For the $128 \times 128$ tile maps used here
this yields a $32 \times 32$ heatmap, i.e.\ $1024$ candidate region centres,
with each heatmap cell covering a $4 \times 4$ tile block. Since a viewport
spans $12 \times 20$ tiles, it occupies $3 \times 5$ cells; the stride is
therefore already coarse relative to the target, and we return to this point in
Section~\ref{sec:discussion}.

Because the heads read that level and no other, the whole backbone is trained,
including the first residual stage. Detection libraries conventionally freeze
the early stages at their ImageNet weights, which is reasonable when a
detection head pools several pyramid levels and scores them with its own
learned layers: the frozen stage is then one input among many. Here it is the
only one, and the input is a channelised tile tensor rather than a natural
image, so leaving it frozen would mean scoring StarCraft game state through
features never trained on anything like it.

\paragraph{Heatmap head.} Two $3{\times}3$ convolutions followed by a
$1{\times}1$ sigmoid layer produce the region-centre heatmap
$\hat{\mathcal{Y}}_t \in [0,1]^{H_s \times W_s}$. The ground-truth heatmap is
rendered by placing a 2D Gaussian at each ranked mode of
Section~\ref{subsec:modes}:
\begin{equation}
  \mathcal{Y}_t(\mathbf{q}) = \max_k\, \omega_k
    \exp\!\left(-\frac{\|\mathbf{q} - \pi(\mathbf{m}^{(k)}_t)\|^2}
                      {2\sigma^2}\right),
  \label{eq:gt_heatmap}
\end{equation}
where the amplitude is
\begin{equation}
  \omega_k = \begin{cases}
    1 & k = 1, \\
    n_t(\mathbf{m}^{(k)}_t)/U \in (0,1] & k \ge 2.
  \end{cases}
  \label{eq:amplitude}
\end{equation}
Using the support as an amplitude for the auxiliary modes, rather than treating
them identically, preserves the ranking in the target and prevents a mode
observed by a single participant from being trained as strongly as one observed
by several.

The primary is exempt. It is already distinguished from the auxiliaries by
which objective supervises it — $\mathcal{L}_{\text{hcm}}$ owns its cell and
$\mathcal{L}_{\text{rmc}}$ owns $\Psi_t$ — so weighting it by support adds no
ranking information and only weakens its supervision. It weakens it worst
where the task is hardest: when observers disagree the Top-1 support is low,
$\omega_1$ is small, and every cell around the peak then carries a
penalty-reduction factor $(1-\omega_1 e^{-\cdot})^{\beta_f}$ close to one and
is trained as a hard negative. The primary would receive a near-delta target on
precisely the multimodal frames this work is about. Ranking across the ordered
set is unaffected, since $\omega_k \le 1$ for $k \ge 2$.

The operator $\pi(\cdot)$ maps a mode to the centre of its assigned heatmap
cell, and it is not cosmetic. A mode at a continuous tile coordinate can fall
up to $s\sqrt{2}/2$ tiles from the nearest sampled grid point, which at
$s = 4$ attenuates the rendered peak substantially. Rendering at the
continuous coordinate would therefore
make the target amplitude encode quantisation error as much as observer
support, defeating the ranking it is meant to carry; a mode with support $1$ of
$U=5$ would render anywhere in $[0.074,\,0.2]$ instead of exactly $0.2$. We
render at the assigned cell, as \cite{zhou2019objects} does, and leave the
sub-cell residual to the offset head.

\paragraph{Offset and size heads.} Following \cite{zhou2019objects},
a sub-pixel offset map $\hat{\mathbf{O}}_t \in \mathbb{R}^{2 \times H_s
\times W_s}$ and a size map $\hat{\mathbf{Z}}_t \in \mathbb{R}^{2 \times
H_s \times W_s}$ recover accurate region boxes from detected centres. Output
convolutions of all three heads are initialised at a small standard deviation,
following \cite{zhou2019objects}; initialising them by the fan-out rule used
for the hidden layers gives a standard deviation of $\sqrt{2}$ for a
$1{\times}1$ convolution with a single output channel, which saturates the
sigmoid and drives the offset head far outside its $[0,1)$ target range, so
that decoded centres fall off the map and are clamped to its border.

\paragraph{Ranked multi-peak extraction.} At inference a $3{\times}3$
max-pooling suppresses non-maxima, and the top-$K$ peaks with score above a
confidence threshold $\tau$ are returned in descending score order. The
highest-scoring peak is the primary region. If fewer than $K$ peaks exceed
$\tau$, fewer regions are returned, so the number of active auxiliary regions
adapts to the number of concurrent events.

Two details of this step are load-bearing. First, $\tau$ must lie below the
smallest meaningful target amplitude $1/U$, since by Eq.~\ref{eq:gt_heatmap} a
mode supported by one observer is trained toward exactly $1/U$; a threshold at
or above that value makes minority regions unreachable no matter how well the
model fits them. Second, \texttt{max\_pool} pads with $-\infty$, so a cell on
the heatmap border competes against fewer neighbours than an interior cell and
becomes a local maximum with probability $1/4$ rather than $1/9$ under a flat
heatmap. We therefore exclude the outermost ring of cells from the candidate
set, which costs nothing reachable: a viewport at the map corner has its centre
at $(h/2, w/2)$, one to two cells inside the boundary.

\paragraph{Differentiable region centres.} The peak indices returned by a
top-$K$ selection are discrete and carry no gradient, so a centre formed as
(peak cell $+$ offset head) is a function of the regression heads alone and is
independent of the heatmap. The repulsion and smoothness objectives below are
statements about \emph{where the model looks}, and defined on such a centre
they cannot reach the map that decides it. We therefore define a soft centre by
a softmax over the heatmap logits in a window of radius $r$ around each peak,
\begin{equation}
  \bar{\mathbf{p}}^{(k)}_t = s\Bigl(\mathbf{c}^{(k)}_t +
    \textstyle\sum_{\boldsymbol{\delta} \in \mathcal{W}_r}
    \mathrm{softmax}\bigl(z_t(\mathbf{c}^{(k)}_t + \boldsymbol{\delta})\bigr)\,
    \boldsymbol{\delta}\Bigr),
  \label{eq:soft_center}
\end{equation}
where $\mathbf{c}^{(k)}_t$ is the $k$-th peak cell, $z_t$ the pre-sigmoid
logits, and $\mathcal{W}_r$ the offsets within the window. This quantity is
differentiable with respect to $\hat{\mathcal{Y}}_t$ and collapses toward the
peak cell as the peak sharpens. It is used only by
$\mathcal{L}_{\text{rep}}$ and $\mathcal{L}_{\text{smooth}}$; the regions
emitted at inference remain the decoded (cell $+$ offset) ones, since adding
both displacements would double-count a sub-cell position that only
$\mathcal{L}_{\text{off}}$ supervises.

\subsection{Multi-Objective Training Losses}
\label{subsec:losses}

Director-CenterNet is trained with four objectives designed for this task,
alongside the standard CenterNet regression losses. The four divide the
heatmap between them: $\mathcal{L}_{\text{hcm}}$ owns the primary cell and the
background, $\mathcal{L}_{\text{rmc}}$ owns the auxiliary modes, and
$\mathcal{L}_{\text{rep}}$ and $\mathcal{L}_{\text{smooth}}$ act on the
geometry of the decoded regions.

\subsubsection{Human Consensus Match Loss $\mathcal{L}_{\text{hcm}}$}

The primary region must be anchored to the Top-1 mode. Writing
$\mathcal{Y}^{(1)}_t$ for the heatmap of Eq.~\ref{eq:gt_heatmap} restricted to
$\mathbf{m}^{(1)}_t$, and $\Psi_t$ for the auxiliary support set defined in
Section~\ref{subsec:rmc_loss}, we apply the modified focal loss
\cite{law2018cornernet}:
\begin{equation}
  \mathcal{L}_{\text{hcm}} = -\frac{1}{N}\sum_{\mathbf{q} \notin \Psi_t}
    \begin{cases}
      (1-\hat{\mathcal{Y}}_t(\mathbf{q}))^{\alpha_f}
        \log \hat{\mathcal{Y}}_t(\mathbf{q})
        & \text{if } \mathbf{q} = \pi(\mathbf{m}^{(1)}_t), \\[4pt]
      (1-\mathcal{Y}_t(\mathbf{q}))^{\beta_f}
        \hat{\mathcal{Y}}_t(\mathbf{q})^{\alpha_f}
        \log(1-\hat{\mathcal{Y}}_t(\mathbf{q}))
        & \text{otherwise,}
    \end{cases}
  \label{eq:hcm_loss}
\end{equation}
where $N$ is the number of positive locations and $\alpha_f, \beta_f$ are
focal loss hyperparameters.

Two aspects of Eq.~\ref{eq:hcm_loss} differ from a direct transcription of
CornerNet and both concern the auxiliary modes. First, the penalty-reduction
factor uses the \emph{joint} target $\mathcal{Y}_t$ rather than
$\mathcal{Y}^{(1)}_t$. In CornerNet the $(1-\mathcal{Y})^{\beta_f}$ factor
protects every ground-truth location from being pulled down as a hard
negative; weighting by $\mathcal{Y}^{(1)}_t$ alone leaves the auxiliary modes
unprotected, since $\mathcal{Y}^{(1)}_t \approx 0$ there and the factor is
therefore $\approx 1$. The focal term would then be trained to erase precisely
what $\mathcal{L}_{\text{rmc}}$ is trying to create. Second, down-weighting by
itself is insufficient: at an auxiliary mode with support $1$ of $U=5$ the
joint target is $0.2$, so $(1-0.2)^{4} = 0.41$ still leaves substantial
downward pressure. We therefore exclude $\Psi_t$ from this loss entirely and
delegate it to $\mathcal{L}_{\text{rmc}}$, so that the two objectives are
complementary by construction rather than by assumption.

The hand-off is conditional on $\mathcal{L}_{\text{rmc}}$ being active, and the
condition matters for reading the ablation. Applied unconditionally, setting
$\lambda_{\text{rmc}} = 0$ removes the positive pull on the auxiliary modes
while still shielding them from suppression, so the ``without
$\mathcal{L}_{\text{rmc}}$'' row is not an ablation of the auxiliary mechanism:
the model places peaks at those locations anyway, because they are genuinely
salient, and nothing pushes them down. Leaving cells both unsupervised and
unsuppressed is not a defensible default either. When $\mathcal{L}_{\text{rmc}}$
is disabled, $\mathcal{L}_{\text{hcm}}$ therefore owns the entire map, which is
the comparison the ablation is meant to draw.

\subsubsection{Ranked Mode Coverage Loss $\mathcal{L}_{\text{rmc}}$}
\label{subsec:rmc_loss}

The auxiliary regions must cover the remaining modes. Let
$\mathcal{M}^-_t = (\mathbf{m}^{(2)}_t, \ldots, \mathbf{m}^{(K)}_t)$ and let
$\mathcal{Y}^-_t$ be the support-weighted heatmap of Eq.~\ref{eq:gt_heatmap}
rendered from $\mathcal{M}^-_t$ alone. Writing $\mathcal{A}^{(1)}_t$ for the
binary mask of tiles already covered by the primary region, the auxiliary
region outside it is
\begin{equation}
  \Omega_t = \bigl\{\mathbf{q} : \mathcal{A}^{(1)}_t(\mathbf{q}) = 0\bigr\},
  \label{eq:omega}
\end{equation}
and the subset of it that actually carries auxiliary mass is
\begin{equation}
  \Psi_t = \bigl\{\mathbf{q} \in \Omega_t : \mathcal{Y}^-_t(\mathbf{q})
    > \epsilon \bigr\}.
  \label{eq:psi}
\end{equation}
The loss is a masked mean-squared error over $\Psi_t$:
\begin{equation}
  \mathcal{L}_{\text{rmc}} = \frac{1}{|\Psi_t|}\sum_{\mathbf{q} \in \Psi_t}
    \bigl(\hat{\mathcal{Y}}_t(\mathbf{q}) - \mathcal{Y}^-_t(\mathbf{q})\bigr)^2.
  \label{eq:rmc_loss}
\end{equation}
When $\mathcal{M}^-_t$ is empty, which occurs in early-game frames where
observers agree, $\mathcal{L}_{\text{rmc}} = 0$ and no auxiliary supervision is
applied.

Normalising over $\Psi_t$ rather than over $\Omega_t$ is essential. $\Omega_t$
excludes only the primary region and therefore covers almost the whole
heatmap — $1009$ of $1024$ cells in our configuration — while the auxiliary
mass occupies roughly a dozen of them. Averaging over $\Omega_t$ dilutes the
only informative cells by two orders of magnitude, whereas
$\mathcal{L}_{\text{hcm}}$ normalises by the number of positives. Suppressing
the true background is already $\mathcal{L}_{\text{hcm}}$'s role; this term
only has to raise the prediction at the ranked minority modes to their
support-weighted amplitude.

\subsubsection{Spatial Repulsion Loss $\mathcal{L}_{\text{rep}}$}

To keep the predicted regions non-redundant we penalise pairwise Intersection
over Union between them, evaluated on the soft centres of
Eq.~\ref{eq:soft_center}:
\begin{equation}
  \mathcal{L}_{\text{rep}} = \sum_{k < l}
    \text{IoU}\!\bigl(V(\bar{\mathbf{p}}^{(k)}_t),\,
                      V(\bar{\mathbf{p}}^{(l)}_t)\bigr).
  \label{eq:rep_loss}
\end{equation}
Regions are shifted rather than cropped at the map boundary, matching the
feasible set $\mathcal{C}$ of Eq.~\ref{eq:feasible}, so that this term always
sees the fixed-size viewport the task is defined over and cannot reduce an
overlap penalty by sliding a region partly off the map.

When fewer than two regions are returned, $\mathcal{L}_{\text{rep}} = 0$. Note
that this term cannot create regions; it only separates regions that
$\mathcal{L}_{\text{rmc}}$ has already produced. Its effect is therefore
defined only conditional on $\mathcal{L}_{\text{rmc}}$, which is why the
corresponding cell of the ablation lattice in Section~\ref{subsec:ablation} is
left empty rather than measured.

\subsubsection{Trajectory Smoothness Loss $\mathcal{L}_{\text{smooth}}$}

To suppress camera thrashing we penalise inter-frame displacement of the
primary region centre. Writing $d_t = \|\bar{\mathbf{p}}^{(1)}_t -
\bar{\mathbf{p}}^{(1)}_{t+1}\|_2$ in tiles and $D = \sqrt{H^2+W^2}$ for the
map diagonal,
\begin{equation}
  \mathcal{L}_{\text{smooth}} = \frac{1}{D\,(T-1)}\sum_{t=1}^{T-1}
    \mathcal{H}_{\Delta_s}(d_t), \qquad
  \mathcal{H}_{\Delta_s}(d) = \begin{cases}
    \tfrac{1}{2}d^2 & d \le \Delta_s, \\[2pt]
    \Delta_s\bigl(d - \tfrac{1}{2}\Delta_s\bigr) & d > \Delta_s.
  \end{cases}
  \label{eq:smooth_loss}
\end{equation}
The objective is applied to the primary region only, since the auxiliary
regions are expected to be reallocated as events appear and disappear.

A squared-$L_2$ penalty in tile units, which is the direct transcription of
the intent, makes ``hold still'' cheaper than ``follow the action''. On a
$128 \times 128$ map two independent predictions differ by
$\mathbb{E}[d^2] \approx 5461$ at initialisation, several orders of magnitude
above a focal term of $O(1)$; because the gradient norm is clipped globally,
such a term does not merely outweigh the others but rescales the whole update
and starves them. The Huber form caps the per-sample gradient at $\Delta_s$ so
that small motion remains almost free while a large jump is expensive but
bounded, and the division by $D$ places the result on an $O(1)$ scale. The
order matters: normalising before applying the Huber would put $\Delta_s$ at
$\Delta_s/D$ and shrink the objective by roughly four orders of magnitude,
disabling it rather than correcting it. We additionally hold
$\lambda_{\text{sm}}$ at zero for the first $E_0$ epochs and ramp it linearly
to its full value by epoch $E_1$, because a trajectory objective imposed on a
randomly initialised heatmap is most cheaply satisfied by not moving at all.

\subsubsection{Joint Training Objective}

The total loss combines the standard CenterNet detection objective with the 
four multi-viewport losses:
\begin{equation}
  \mathcal{L} = \mathcal{L}_{\text{hcm}} 
              + \lambda_{\text{off}}\mathcal{L}_{\text{off}} 
              + \lambda_{\text{sz}}\mathcal{L}_{\text{size}} 
              + \lambda_{\text{rmc}}\mathcal{L}_{\text{rmc}}
              + \lambda_{\text{rep}}\mathcal{L}_{\text{rep}}
              + \lambda_{\text{sm}}\mathcal{L}_{\text{smooth}},
  \label{eq:total_loss}
\end{equation}
where $\mathcal{L}_{\text{off}}$ and $\mathcal{L}_{\text{size}}$ are the 
standard CenterNet offset and size regression losses \cite{zhou2019objects}.

Table~\ref{tab:loss_summary} summarises the role of each loss component.

\begin{table}[t]
  \centering
  \caption{Summary of Director-CenterNet training losses.}
  \label{tab:loss_summary}
  \begin{tabular}{llll}
    \toprule
    Loss & Symbol & Type & Purpose \\
    \midrule
    Human Consensus Match & $\mathcal{L}_{\text{hcm}}$ & Focal & 
      Anchor main viewport to human consensus \\
    Offset regression & $\mathcal{L}_{\text{off}}$ & L1 & 
      Sub-pixel centre accuracy \\
    Size regression & $\mathcal{L}_{\text{size}}$ & L1 & 
      Viewport size accuracy \\
    Ranked Mode Coverage & $\mathcal{L}_{\text{rmc}}$ & Masked MSE &
      Direct auxiliary regions to the remaining ranked modes \\
    Spatial Repulsion & $\mathcal{L}_{\text{rep}}$ & IoU penalty &
      Keep the predicted regions non-redundant \\
    Trajectory Smoothness & $\mathcal{L}_{\text{smooth}}$ & Huber &
      Suppress thrashing of the primary region \\
    \bottomrule
  \end{tabular}
\end{table}

%%====================================================================
\section{Experiments}
\label{sec:experiments}
%%====================================================================

\subsection{Dataset and Preprocessing}

We use the same StarCraft: Brood War dataset and Observer preprocessing 
pipeline as \cite{joo2023learning, anonymous2025kbrs}. Replay files were 
collected from online communities (Ygosu, BWReplays) with constraints on 
player APM ($>250$), race matchups (Z vs.~P, T vs.~P, Z vs.~T), competitive 
maps (Fighting Spirit, Jade, Othello), and starting locations. Human
observational data were collected from participants (aged 20--35) who freely
spectated replay files while viewport coordinates were recorded at each frame.
A total of approximately 420{,}000 frames from 21 games were collected.

\paragraph{Observer pool.} We use the observational data of
\cite{joo2023learning} unchanged, with no additional collection: $U=5$
observers per replay, the same replay files, and the same three replay-level
folds. Reusing it without modification is deliberate. The contribution of this
work is a different reading of that data — the minority modes previous work
suppressed — so holding the data fixed makes the comparison against the
previously reported single-region results direct rather than confounded by a
change of corpus. It also bounds the ranking at five levels of support, which
Section~\ref{sec:discussion} discusses.

\subsection{Structure of the Target Distribution}
\label{subsec:target_structure}

Before evaluating any proposed model we characterise the supervision itself.
The mode counts below are a property of the aggregated observer data alone; the
accuracy columns are measured with Mask R-CNN + prior \cite{anonymous2025kbrs},
the strongest previously reported single-region observer, over $148{,}638$ test
frames. This analysis is a contribution in its own right, since it quantifies
the multimodality that motivates the task and, as shown below, identifies the
variable that actually governs single-region accuracy.

\begin{table}[t]
  \centering
  \caption{Single-region accuracy as a function of the number of modes in the
  aggregated observer distribution. IR falls monotonically as the target
  becomes more multimodal, and the fraction of frames whose predicted region
  serves a minority mode rather than the Top-1 mode rises correspondingly.
  Measured with Mask R-CNN + prior over $148{,}638$ frames.}
  \label{tab:by_nmodes}
  \begin{tabular}{lrrccc}
    \toprule
    $|\mathcal{M}_t|$ & Frames & Share & IR $\uparrow$ & I@0.5 $\uparrow$
      & Served by minority mode \\
    \midrule
    1 & 61{,}397 & 41.3\% & 0.779 & 0.830 & 0.000 \\
    2 & 19{,}702 & 13.3\% & 0.714 & 0.774 & 0.256 \\
    3 & 33{,}962 & 22.8\% & 0.609 & 0.682 & 0.284 \\
    4 & 28{,}242 & 19.0\% & 0.546 & 0.626 & 0.337 \\
    5 &  5{,}335 &  3.6\% & 0.499 & 0.577 & 0.479 \\
    \bottomrule
  \end{tabular}
\end{table}

Table~\ref{tab:by_nmodes} shows that $58.7\%$ of frames contain at least two
modes, and that IR degrades by $36\%$ relative between unimodal frames and
five-mode frames. The final column is the more direct statement: as the target
fragments, a model trained to produce one region increasingly produces a region
that serves a \emph{minority} of observers rather than the consensus. It is
already behaving multimodally, but without control over which mode it selects.

\begin{table}[t]
  \centering
  \caption{Single-region accuracy and temporal behaviour as a function of the
  support margin $n_t(\mathbf{m}^{(1)}_t) - n_t(\mathbf{m}^{(2)}_t)$. The final
  column is the fraction of consecutive frame pairs on which the prediction
  alternates between the top two modes. Measured with Mask R-CNN + prior.}
  \label{tab:by_margin}
  \begin{tabular}{lrcccc}
    \toprule
    Margin & Frames & IR $\uparrow$ & I@0.5 $\uparrow$
      & Served by minority mode & Top-2 alternation \\
    \midrule
    0 & 26{,}917 & 0.595 & 0.670 & 0.360 & 0.0706 \\
    1 & 47{,}676 & 0.609 & 0.681 & 0.303 & 0.0413 \\
    2 & 12{,}732 & 0.618 & 0.688 & 0.204 & 0.0383 \\
    3 & 17{,}727 & 0.709 & 0.768 & 0.000 & 0.0109 \\
    4 & 27{,}491 & 0.780 & 0.832 & 0.000 & 0.0023 \\
    5 & 16{,}095 & 0.853 & 0.895 & 0.000 & 0.0000 \\
    \bottomrule
  \end{tabular}
\end{table}

Table~\ref{tab:by_margin} isolates the thrashing mechanism. When the top two
modes are supported by equal numbers of observers, the single-region model
alternates between them on $7.1\%$ of consecutive frame pairs; when one mode is
unanimous, it never does. Camera thrashing is therefore not a generic
instability of the model but a specific consequence of forcing a single output
on an ambiguous target, and it is worst exactly where the ambiguity is
greatest.

\begin{table}[t]
  \centering
  \caption{Target structure and single-region accuracy by game progression.
  Mode counts are a property of the observer data; the remaining columns are
  measured with Mask R-CNN + prior. Accuracy is \emph{not} monotonic in
  progression: it tracks the mode count, which peaks at mid-game and falls as
  the match resolves.}
  \label{tab:by_progression}
  \begin{tabular}{lrcccc}
    \toprule
    Progression & Frames & $|\mathcal{M}_t|$ & IR $\uparrow$
      & OC3@0.5 $\uparrow$ & VD $\downarrow$ \\
    \midrule
    1/4 & 37{,}163 & 2.329 & 0.734 & 0.661 & 3.56 \\
    2/4 & 37{,}157 & 2.596 & 0.572 & 0.509 & 4.81 \\
    3/4 & 37{,}158 & 2.422 & 0.654 & 0.525 & 3.61 \\
    4/4 & 37{,}160 & 1.865 & 0.749 & 0.632 & 1.72 \\
    \bottomrule
  \end{tabular}
\end{table}

Table~\ref{tab:by_progression} qualifies a pattern reported in earlier work.
Previous studies describe accuracy as declining with game progression
\cite{joo2023learning, anonymous2025kbrs}; on our test set the relation is not
monotonic but U-shaped, with the minimum at the midpoint and the maximum at the
end. The mode count follows the inverse shape, and the rank ordering of the two
columns agrees at every bin. We therefore read progression as a proxy for
multimodality rather than as the governing variable: mid-game is where events
occur at several locations at once, whereas a late game that has resolved into
one decisive engagement is easier to observe than an early one. VD behaves the
same way, which is consistent with thrashing being driven by ambiguity rather
than by pace.

Game-state tensors follow the channelisation protocol of \cite{bae2025observer}:
four entity channels per player (workers, ground units, air units, buildings),
terrain elevation, combined vision, and resource channels. Temporal stacking
uses $\Delta=8$, $W=3$ (four-frame stack), and the $128 \times 128$ tile map
is consumed at its native resolution.

\subsection{Experimental Setup}

\paragraph{Splits.} Three replay-level folds are used for cross-validation 
following \cite{joo2023learning, anonymous2025kbrs}. All main comparison 
results are averaged over three folds and three random seeds 
$\{123, 456, 789\}$.

\paragraph{Implementation.} Backbone: ResNet-50 pretrained on ImageNet with
FPN, finest level ($s=4$), giving a $32 \times 32$ heatmap, all stages
trainable. Heads are two $3{\times}3$ convolutions of width $64$ followed by a
$1{\times}1$ output layer. Gaussian rendering: $\sigma = 4$~tiles, one heatmap
cell at this stride; at $\sigma = 2$ the target is half a cell wide, so the
immediate neighbours of a peak fall to $e^{-2} = 0.135$ and are trained as
near-hard negatives, leaving no soft neighbourhood at all. Mode extraction:
Gaussian smoothing
$\sigma_m = 4$~tiles, relative floor $\theta = 0.35$, minimum separation
$D = 12$~tiles. Loss weights: $\lambda_{\text{off}}=1$,
$\lambda_{\text{sz}}=0.1$, $\lambda_{\text{rmc}}=0.5$,
$\lambda_{\text{rep}}=0.3$, $\lambda_{\text{sm}}=0.1$. Smoothness shaping:
Huber knee $\Delta_s = 5$~tiles, warmup $E_0 = 5$, $E_1 = 10$. Soft-centre
window radius $r=2$~cells. Region set: $K=3$, confidence threshold
$\tau=0.1$. Note that $\tau$ is set below $1/U$ for the reasons given in
Section~\ref{sec:method}; the inference-time score filter is held at the same
value, since a filter above $\tau$ would silently remove every minority
region.

\paragraph{Optimiser.} SGD with momentum~0.9, weight decay $5{\times}10^{-4}$, 
learning rate $5{\times}10^{-3}$ with StepLR decay (step~3, factor~0.1), 
batch size~16, 30~epochs.

\subsection{Baselines}

We compare Director-CenterNet against the following methods.

\begin{itemize}
  \item \textbf{BC}: Behaviour cloning via MSE regression to human viewport 
        coordinates.
  \item \textbf{AIIDE}: Rule-based observer \cite{mattsson2015automatic}.
  \item \textbf{SSCAIT}: Rule-based observer with event priorities and timers 
        \cite{mattsson2015automatic}.
  \item \textbf{Mask R-CNN (1H)}: Single-viewport detection, one human label 
        \cite{joo2023learning}.
  \item \textbf{Mask R-CNN (5H)}: Single-viewport detection, five human labels 
        \cite{joo2023learning}.
  \item \textbf{Mask R-CNN + ROCI}: With Region of Common Interest loss 
        \cite{joo2023learning}.
  \item \textbf{Mask R-CNN + fixed-kernel prior}: the auxiliary-prior variant
        of \cite{anonymous2025kbrs}, the strongest previously reported
        single-region result and our primary point of comparison.
  \item \textbf{CenterNet (single)}: the proposed backbone trained with
        $\mathcal{L}_{\text{hcm}}$ only, i.e.\ the first row of
        Table~\ref{tab:ablation}. This isolates the effect of the backbone
        change from that of the multi-region formulation, and is therefore
        the same architecture as Director-CenterNet rather than a separate
        plain-backbone CenterNet, which would confound the two.
\end{itemize}

\subsection{Main Results}

Table~\ref{tab:main_results} reports main-viewport performance (IR and 
Intersection@$\delta$) on the standard test set.

\begin{table}[t]
  \centering
  \caption{Main viewport performance on the standard test set (fixed race, 
  map, starting location). Results averaged over three folds $\times$ three 
  seeds. Best result in \textbf{bold}, second best \underline{underlined}.
  The two missing rows require the full three-fold, three-seed protocol and
  the @any and @0.3 thresholds; fold-1, seed-456 values at @0.5 are in
  Table~\ref{tab:multi_results} and are not substitutable here.
  \todo{Fill in experimental values.}}
  \label{tab:main_results}
  \begin{tabular}{lcccc}
    \toprule
    Method & @any $\uparrow$ & @0.3 $\uparrow$ & @0.5 $\uparrow$ 
           & IR $\uparrow$ \\
    \midrule
    BC                         & 0.061 & 0.059 & 0.046 & -- \\
    AIIDE                      & 0.524 & --    & --    & -- \\
    SSCAIT                     & 0.491 & --    & --    & -- \\
    Mask R-CNN (1H)            & 0.593 & --    & --    & -- \\
    Mask R-CNN (5H)            & 0.569 & --    & --    & -- \\
    Mask R-CNN + ROCI          & 0.588 & --    & --    & -- \\
    Mask R-CNN + prior         & \underline{0.824} & \underline{0.798}
                               & \underline{0.770} & \underline{0.709} \\
    CenterNet (single)         & \todo{X} & \todo{X} & \todo{X} & \todo{X} \\
    \midrule
    \textbf{Director-CenterNet (Ours)} & \todo{X} & \todo{X} 
                               & \todo{X} & \todo{X} \\
    \bottomrule
  \end{tabular}
\end{table}

Table~\ref{tab:multi_results} reports the multi-region metrics. These are the
quantities the proposed formulation is designed to improve: the gap between
I@$\delta$ and BoK@$\delta$ indicates how much of the human attention
distribution is reachable by additional regions, and OC@$\delta$ indicates how
many of the disagreeing observers the output actually serves. Every method is
scored on its own top-$K$ regions, including the baselines, which as noted in
Section~\ref{sec:problem} emit a ranked list rather than a single box.

\begin{table}[t]
  \centering
  \caption{Multi-region coverage, per-observer service, and temporal stability
  on the mode-disagreement evaluation set ($\approx 148{,}600$ frames, fold~1
  test replays, seed $456$ throughout). $|\hat{\mathcal{P}}_t|$ is the mean
  number of regions the model emits per frame; BoK and OC are evaluated on its
  top $K$ of those. Every method answers on $\ge 99.9\%$ of frames at its own
  default confidence threshold, so the values are directly comparable without a
  coverage correction. Lower VD and jerk are better. \todo{The CenterNet
  (single) row is still the seed-123 measurement and is not comparable to the
  rest; re-run it at seed 456. Then re-measure all rows under the three-fold,
  three-seed protocol of Table~\ref{tab:main_results}.}}
  \label{tab:multi_results}
  \begin{tabular}{lrccccrr}
    \toprule
    & & \multicolumn{2}{c}{$K=1$} & \multicolumn{2}{c}{$K=3$} & & \\
    \cmidrule(lr){3-4}\cmidrule(lr){5-6}
    Method & $|\hat{\mathcal{P}}_t|$ & BoK@0.5 & OC@0.5 & BoK@0.5 & OC@0.5
           & VD $\downarrow$ & Jerk $\downarrow$ \\
    \midrule
    Mask R-CNN            & 6.99 & 0.709 & 0.343 & 0.914 & 0.570 & 4.43 & 16.4 \\
    Mask R-CNN + prior    & 6.82 & \textbf{0.744} & \textbf{0.371}
                          & \textbf{0.929} & \textbf{0.592}
                          & \textbf{3.19} & 12.5 \\
    CenterNet (single)$^{\dagger}$ & 4.74 & 0.680 & 0.328 & 0.915 & 0.557
                          & 7.20 & 27.8 \\
    \midrule
    \textbf{Director-CenterNet} & 2.72 & 0.668 & 0.313 & 0.896 & 0.516
                          & 3.30 & \textbf{12.3} \\
    \bottomrule
  \end{tabular}
  \\[2pt]
  \footnotesize $^{\dagger}$Seed 123; see \todo{above}.
\end{table}

Three observations follow, and together they set out what the proposed method
does and does not establish.

First, the baselines are not single-region systems. They emit $4.7$ to $7.0$
boxes per frame, with at least three in over $93\%$ of frames; previous work
simply reports the first. Scored on their top three they reach BoK3@0.5 of
$0.91$--$0.93$, so almost all of the human attention union is already reachable
from an untrained ranked list and there is little headroom in that metric. The
informative quantity is OC@$\delta$, which asks how many \emph{individual}
observers are served: it rises only from $0.371$ to $0.592$ as $K$ goes from
$1$ to $3$, and stops well short of $1$. The additional boxes a detector
produces are proposals clustered around the same evidence, not a ranked account
of where the disagreeing observers were looking.

Second, the backbone change is not free. Replacing Mask R-CNN with a CenterNet
heatmap detector costs a factor of two in temporal stability (VD $3.19
\rightarrow 7.20$, jerk $12.5 \rightarrow 27.8$) before any multi-region
objective is introduced. The Trajectory Smoothness loss is therefore not a
refinement but a requirement of the architecture we adopt, and the appropriate
reference point for it is the Mask R-CNN baseline rather than the untreated
CenterNet. Measured that way it succeeds: the full model recovers jerk to
$12.3$ and VD to $3.30$, restoring the stability the architecture gave up.

Third — and this is the result that determines how the rest of the paper should
be read — Director-CenterNet does \emph{not} lead this table. Mask R-CNN +
prior is ahead on every accuracy column, by $0.08$ in IR and $0.08$ in OC3@0.5.
What separates the two methods is the last column that is not shown here:
Mask R-CNN reaches those numbers while emitting $6.8$ regions per frame against
Director-CenterNet's $2.7$. A comparison at that ratio does not measure
multi-region quality; it measures how many guesses each method is permitted.
Section~\ref{subsec:budget} therefore constrains the baselines to a comparable
region budget and asks what survives.

\subsection{Region Count Under a Matched Budget}
\label{subsec:budget}

A detector emits as many regions as its confidence threshold admits, so the
comparison in Table~\ref{tab:multi_results} confounds two things: how good each
region is, and how many the method was allowed to produce. This section
separates them by sweeping the baselines' score threshold over $\{0.5, 0.7,
0.9\}$, which is the only mechanism a trained detector offers for controlling
its region count, and reporting the same metrics at each operating point
(Table~\ref{tab:threshold}).

\begin{table}[t]
  \centering
  \caption{Baseline score-threshold sweep against Director-CenterNet
  (fold~1, seed~$456$, $148{,}644$ evaluated frames). Cov.\ is frame coverage;
  IR$^{\ast}$ and OC3$^{\ast}$ are coverage-penalised, scoring an unanswered
  frame as $0$. $\beta_n$ is the count responsiveness of
  Eq.~\eqref{eq:count_slope} and $\text{MAE}_n$, $\text{Exact}_n$ the count
  accuracy of Eq.~\eqref{eq:count_err}. The shaded row is the operating point
  whose region budget matches ours.}
  \label{tab:threshold}
  \begin{tabular}{llrrcccrrr}
    \toprule
    Method & $\tau$ & Cov.\ & $|\hat{\mathcal{P}}_t|$
           & IR & IR$^{\ast}$ & OC3$^{\ast}$
           & $\beta_n$ & $\text{MAE}_n$ & $\text{Exact}_n$ \\
    \midrule
    Mask R-CNN            & 0.5 & 100.0\% & 6.99 & 0.641 & 0.641 & 0.569
                          & $-0.198$ & 4.76 & \phantom{0}4.0\% \\
                          & 0.7 & \phantom{0}92.2\% & 2.64 & 0.623 & 0.574 & 0.439
                          & $-0.148$ & 1.62 & 18.4\% \\
                          & 0.9 & \phantom{0}55.7\% & 1.42 & 0.717 & 0.399 & 0.268
                          & $-0.043$ & 1.05 & 39.1\% \\
    \midrule
    Mask R-CNN            & 0.5 & \phantom{0}99.9\% & 6.82 & 0.684
                          & \textbf{0.683} & \textbf{0.592}
                          & $-0.279$ & 4.65 & \phantom{0}4.5\% \\
    \quad + prior         & 0.7 & \phantom{0}90.2\% & 2.76 & 0.664 & 0.599 & 0.469
                          & $-0.143$ & 1.63 & 17.9\% \\
                          & 0.9 & \phantom{0}54.8\% & 1.47 & 0.762 & 0.418 & 0.294
                          & $-0.052$ & \textbf{0.96} & \textbf{40.7\%} \\
    \midrule
    \textbf{Ours}         & 0.1 & 100.0\% & 2.72 & 0.601 & 0.601 & 0.516
                          & $\mathbf{+0.135}$ & 1.04 & 29.6\% \\
    \bottomrule
  \end{tabular}
\end{table}

\paragraph{Raising the threshold does not improve the prediction; it deletes
frames.} The IR column appears to rise at $\tau=0.9$, to $0.762$ for
Mask R-CNN + prior, above every other entry in the table. It is an artefact of
conditioning. At that threshold the model answers on $54.8\%$ of frames, and
the frames it declines are systematically the hard ones: $60\%$ of them have
three or four ranked modes, against $42\%$ in the evaluation set as a whole.
Restricted to the $64{,}561$ frames that every method in the table answers,
$\tau=0.7$ and $\tau=0.9$ give IR of $0.7753$ and $0.7754$ and
I@0.5 of $0.8900$ and $0.8900$ — identical to four decimal places. Thresholding
cannot improve the primary region, because the primary region is the
highest-scoring box and raising the bar never changes which box that is. It
only removes lower-ranked boxes, and eventually the frame. The
coverage-penalised columns report what remains once that is accounted for.

\paragraph{Count \emph{level} is a threshold; count \emph{direction} is not.}
$\text{MAE}_n$ falls from $4.65$ to $0.96$ and $\text{Exact}_n$ rises from
$4.5\%$ to $40.7\%$ as the threshold is raised, ending better than
Director-CenterNet on both. This is expected and it is worth stating plainly:
the number of emitted regions is a free parameter, and moving its mean onto the
mean of $n_t$ is a calibration exercise that requires no knowledge of the scene.
The responsiveness $\beta_n$ behaves differently. It is negative for every
baseline at every threshold — $-0.279$, $-0.143$, $-0.052$ for the prior
variant — and approaches zero from below without changing sign. A negative
$\beta_n$ means the detector emits \emph{fewer} regions on frames whose target
has more modes, which is the opposite of the required behaviour. The mechanism
is visible in the architecture: a proposal score is a per-box confidence, and in
a frame whose evidence is genuinely divided that confidence is spread across
more boxes, so fewer of them clear a fixed bar. Thresholding can move the level
of the count but not its dependence on the scene, because that dependence is a
property of how the scores are formed. Director-CenterNet attains
$\beta_n = +0.135$.

The sign belongs to the representation and the decoding, not to the objectives
proposed here. Measured across the ablation lattice, $\beta_n$ is $+0.151$ with
$\mathcal{L}_{\text{hcm}}$ alone and between $+0.135$ and $+0.155$ for every
other combination; adding the ranked-mode and repulsion objectives does not
raise it. What produces the positive sign is scoring each peak against its own
neighbourhood and admitting it on its own merits, so a frame with three
well-supported modes yields three peaks that each clear the threshold
independently. We state this explicitly because the natural reading — that
supervising auxiliary modes is what teaches the model to emit more of them — is
not what the data show.

\paragraph{At a matched budget the ranking reverses.} The row to compare
against is $\tau = 0.7$, where Mask R-CNN + prior emits $2.76$ regions per
frame against our $2.72$. There, Director-CenterNet is ahead on
coverage-penalised IR ($0.601$ vs.\ $0.599$), on coverage-penalised OC3@0.5
($0.516$ vs.\ $0.469$), and on all three count statistics
($\beta_n$ $+0.135$ vs.\ $-0.143$; $\text{MAE}_n$ $1.04$ vs.\ $1.63$;
$\text{Exact}_n$ $29.6\%$ vs.\ $17.9\%$). Mask R-CNN's advantage in
Table~\ref{tab:multi_results} is therefore purchased with region count: it is
ahead when permitted $2.5\times$ as many guesses, and behind when constrained to
the same number. We do not read this as the proposed method dominating —
$\tau = 0.5$ remains the better operating point in absolute terms if one is
willing to display seven regions, and Section~\ref{subsec:limitations} states
what that means for the claim. We read it as establishing that the region
budget, not the region quality, is what the headline comparison was measuring.

\subsection{Ablation Study}
\label{subsec:ablation}

\paragraph{Design of the lattice.} Four objectives are designed for this task,
but only three of them are ablation axes. $\mathcal{L}_{\text{off}}$ and
$\mathcal{L}_{\text{size}}$ are the standard CenterNet regression losses and
are a precondition for decoding a box at all, not a contribution whose effect
is measured. $\mathcal{L}_{\text{hcm}}$ is the base rather than an axis:
setting $\lambda_{\text{hcm}}=0$ leaves no positive supervision anywhere on the
heatmap, because $\mathcal{L}_{\text{rmc}}$'s domain is restricted to $\Psi_t$
and supervises neither the primary cell nor the background. The lattice is
therefore $2^3 = 8$ rather than $2^4 = 16$.

Two of the eight cells are not measured, and for different reasons.
$\mathcal{L}_{\text{hcm}} + \mathcal{L}_{\text{rep}}$ is undefined by
construction: as noted in Section~\ref{subsec:losses},
$\mathcal{L}_{\text{rep}}$ penalises overlap between regions that already
exist and cannot create any, so without $\mathcal{L}_{\text{rmc}}$ there is
effectively one region per frame and nothing to repel. The effect of
$\mathcal{L}_{\text{rep}}$ is thus defined only conditional on
$\mathcal{L}_{\text{rmc}}$, and the empty cell records that dependency rather
than an omission. $\mathcal{L}_{\text{hcm}} + \mathcal{L}_{\text{smooth}}$ is
omitted as redundant: the effect of the smoothness objective is read directly
from VD between the third and fourth rows of Table~\ref{tab:ablation}.

Table~\ref{tab:ablation} reports the cumulative series and
Table~\ref{tab:ablation_loo} the leave-one-out series; the third row of the
former and the last row of the latter are the same configuration.

\begin{table}[t]
  \centering
  \caption{Cumulative ablation (fold 1, seed 456, $148{,}644$ frames,
  full-length evaluation). \checkmark~= included. OC is OC3@0.5. No row differs
  from any other by more than $0.015$ in IR, $0.017$ in OC3 or $0.42$ in VD,
  against measured seed-to-seed differences of $0.021$, $0.015$ and $0.38$ for
  one fixed configuration. The entries are reported for completeness; none of
  them is separable from noise at one seed, VD included. \todo{Re-run at seeds
  123 and 789 and report the standard deviation; until then no row ordering
  here should be claimed.}}
  \label{tab:ablation}
  \begin{tabular}{cccccccc}
    \toprule
    $\mathcal{L}_{\text{hcm}}$ &
    $\mathcal{L}_{\text{rmc}}$ &
    $\mathcal{L}_{\text{rep}}$ &
    $\mathcal{L}_{\text{smooth}}$ &
    IR $\uparrow$ & OC@0.5 $\uparrow$ & $|\hat{\mathcal{P}}_t|$ & VD $\downarrow$ \\
    \midrule
    \checkmark &            &            &            & 0.607 & 0.527 & 2.71 & 3.72 \\
    \checkmark & \checkmark &            &            & 0.607 & 0.528 & 2.71 & 3.58 \\
    \checkmark & \checkmark & \checkmark &            & 0.592 & 0.523 & 2.72 & 3.45 \\
    \checkmark & \checkmark & \checkmark & \checkmark & 0.601 & 0.516 & 2.72 & 3.30 \\
    \bottomrule
  \end{tabular}
\end{table}

\begin{table}[t]
  \centering
  \caption{Leave-one-out ablation (fold 1, seed 456). The
  $\mathcal{L}_{\text{hcm}} + \mathcal{L}_{\text{rep}}$ cell of the lattice is
  omitted because $\mathcal{L}_{\text{rep}}$ has no effect without
  $\mathcal{L}_{\text{rmc}}$. The same caveat as Table~\ref{tab:ablation}
  applies, to every column: no row is separated from another by more than the
  measured seed-to-seed variation.}
  \label{tab:ablation_loo}
  \begin{tabular}{lcccc}
    \toprule
    Configuration & IR $\uparrow$ & OC@0.5 $\uparrow$ & $|\hat{\mathcal{P}}_t|$ & VD $\downarrow$ \\
    \midrule
    Full                                    & 0.601 & 0.516 & 2.72 & \textbf{3.30} \\
    w/o $\mathcal{L}_{\text{rmc}}$          & 0.598 & 0.525 & 2.74 & 3.39 \\
    w/o $\mathcal{L}_{\text{rep}}$          & 0.607 & 0.528 & 2.71 & 3.56 \\
    w/o $\mathcal{L}_{\text{smooth}}$       & 0.592 & 0.523 & 2.72 & 3.45 \\
    \bottomrule
  \end{tabular}
\end{table}

\paragraph{What the lattice does and does not show.} Read across either table,
the three optional objectives are not separable on the accuracy metrics at one
seed. The full span of IR across all seven measured configurations is $0.592$
to $0.618$, and the seed-to-seed standard deviation of a single configuration
is of the same order. We therefore make no claim about the relative
contribution of $\mathcal{L}_{\text{rmc}}$ and $\mathcal{L}_{\text{rep}}$ from
these data.

Nor does temporal stability separate them, although at first it appeared to.
Adding $\mathcal{L}_{\text{smooth}}$ moves VD from $3.45$ to $3.30$ and jerk
from $13.4$ to $12.3$, and both are in the expected direction; but a replicate
of the full configuration at a second seed moves VD by $0.38$, which is more
than twice the effect. The direction is consistent across every row that
carries the objective and the mechanism is not in doubt, but on one seed the
magnitude is not measured. \todo{Three seeds per configuration. Until then this
table supports no ordering on any metric, and the accompanying text should not
be read as one.}

An earlier formulation of $\mathcal{L}_{\text{smooth}}$, penalising squared
displacement without normalisation, is worth recording because it failed in a
way that is easy to reproduce. At $\lambda_{\text{sm}} = 0.2$ every
configuration carrying the term collapsed to IR $\approx 0.035$: with
displacements measured in tiles on a $128 \times 128$ map, the squared penalty
reached four orders of magnitude above the other objectives, and under a single
global gradient-norm clip it consumed the entire gradient budget. The cheapest
way to satisfy it is to stop moving, which the model duly learned. The Huber
form of Eq.~\eqref{eq:smooth_loss}, normalised by the map diagonal, is what makes
the term usable, and the warmup exists because the objective is a statement
about a camera that is already tracking something.

\paragraph{Reporting.} $|\hat{\mathcal{P}}_t|$ denotes the mean number of
regions emitted per frame, which we report alongside the accuracy metrics
because $\mathcal{L}_{\text{rmc}}$ and $\mathcal{L}_{\text{rep}}$ are not
interpretable when it is close to one: with a single region per frame, OC@$\delta$
degenerates to the single-region case and $\mathcal{L}_{\text{rep}}$ has no
argument. We also report the seed standard deviation, since several of the
differences between neighbouring rows are of the same order as run-to-run
variation. \todo{Report seed std once the three seeds complete.}

\subsection{Qualitative Analysis}

\todo{Insert qualitative figures:
(1) Side-by-side: single-viewport vs.\ Director-CenterNet main + PIP during 
    concurrent game events (battle + resource harassment). 
(2) Heatmap visualisation: aggregated observer coverage $a_t$, its smoothed
    field with the extracted ranked modes marked, the predicted heatmap
    $\hat{\mathcal{Y}}_t$, and the extracted region centres.
(3) Temporal trajectory plots: frame-by-frame main viewport centre position 
    for Mask R-CNN + prior vs.\ Director-CenterNet, illustrating reduced VD.}

%%====================================================================
\section{Discussion}
\label{sec:discussion}
%%====================================================================

\subsection{Why Multi-Region Finding Matters for Esports Broadcasting}

The fundamental tension in Esports broadcasting is that a single camera 
cannot cover the full extent of a large, concurrent game state. Real-time 
strategy games such as StarCraft exhibit up to 40 times the unit count of 
MOBA games \cite{joo2023learning}, making multi-region coverage not merely 
a convenience but a broadcast necessity. Director-CenterNet addresses this
need directly, and does so using supervision that was already present in the
data prior work collected: the minority modes of the observer distribution
that ROCI and fixed-kernel priors were designed to suppress. No new form of
annotation is introduced, only a denser sampling of the same protocol.

It is worth being precise about what is and is not new here, because
Table~\ref{tab:multi_results} shows that the baselines already emit several
regions per frame. Producing more than one region is not the contribution; a
detector does that by construction, and previous work simply reported only the
first. What a detector does not do is decide \emph{how many} regions the
current frame warrants, order them by how much of the audience each represents,
or keep them from duplicating one another. Those three properties are what turn
a proposal list into a broadcast layout, and they are what the ranked-mode
supervision, the repulsion penalty and the confidence threshold supply. The
measurable consequence is OC@$\delta$: the baselines' additional proposals
raise per-observer service only from $0.371$ to $0.592$, because they cluster
around the same evidence rather than covering the observers who disagreed.

\subsection{CenterNet vs.\ Mask R-CNN for MRVP}

The argument for a heatmap detector in this setting is usually made
structurally: a dense heatmap retains the multimodal structure of the target
where an instance mask forces a discrete assignment that discards it,
local-maximum extraction yields a variable number of regions without changing
the model, and support-weighted Gaussian rendering encodes the ranking of human
attention rather than merely its location. Our results let one of those claims
be tested rather than asserted, and the test is sharper than the structural
statement.

The testable claim is that the two architectures differ in how the emitted
region count depends on the scene. Section~\ref{subsec:budget} measures this as
$\beta_n$ and finds a sign difference that no confidence threshold removes:
every Mask R-CNN operating point is negative ($-0.279$ to $-0.043$),
Director-CenterNet is positive ($+0.135$). The reason is specific to how the
two score their outputs. A proposal score is a per-box confidence competing
against other boxes for the same evidence; when the evidence is genuinely
divided across several regions, each box's share of it falls, so raising a fixed
threshold prunes exactly the frames that needed several regions. A heatmap peak
is scored against its own neighbourhood, so a frame with three well-supported
modes produces three peaks that each clear the threshold on their own merits.
This is why thresholding moves the level of $|\hat{\mathcal{P}}_t|$ but never
its sign: the level is a decision rule, the sign is a property of the
representation.

Two qualifications keep this from being a general claim of superiority. First,
the property belongs to the representation and the decoding rather than to the
objectives this paper proposes: across the ablation lattice $\beta_n$ stays
between $+0.135$ and $+0.155$, and it is highest with
$\mathcal{L}_{\text{hcm}}$ alone. The ranked-mode supervision is what makes the
auxiliary regions correspond to \emph{which} observers disagreed; it is not
what makes the model emit more of them. Second, the architecture change is costly
elsewhere. Untreated CenterNet is twice as unstable as Mask R-CNN
(Table~\ref{tab:multi_results}), which is why $\mathcal{L}_{\text{smooth}}$ is
a requirement rather than a refinement, and at its best operating point
Mask R-CNN + prior remains ahead on IR. The honest summary is that the heatmap
formulation buys a qualitatively different count behaviour and pays for it in
single-region accuracy and in stability that must be restored by an additional
objective.

\subsection{Heatmap Resolution}

The finest FPN level gives a stride of $4$, so a $12 \times 20$-tile viewport
spans only $3 \times 5$ heatmap cells. Relative to target size this is
considerably coarser than the regime CenterNet was designed for, where objects
span tens of cells. Positional precision is not the binding constraint, since
the offset head decodes sub-cell position continuously; what a coarse stride
limits is the density of candidate centres and the separability of nearby
modes, both of which bear on the multi-region metrics rather than on IR.

Resolution is consequently the first place to look for the IR deficit in
Table~\ref{tab:multi_results}, and the comparison is not a fair one until it
has been ruled out: Mask R-CNN pools its features at ROI resolution and is not
subject to the same constraint. A study of this is in progress but not
reported here. We note two requirements that a correct version of it must
satisfy, because both are easy to violate silently. Several constants in the
method are distances in tiles that happen to be expressed in heatmap cells —
the non-maximum suppression radius, which must remain commensurate with the
minimum mode separation $D$; the soft-argmax window; and the excluded border
ring — and each shrinks with the stride unless it is re-derived, the first of
these turning a single mode into several duplicate peaks inside its own
viewport. Separately, $\mathcal{L}_{\text{hcm}}$ sums its negative term over
cells while the normaliser $N$ of Eq.~\ref{eq:hcm_loss} is $1$ per frame by
construction — there is exactly one positive location,
$\pi(\mathbf{m}^{(1)}_t)$ — so its magnitude grows
with the square of the resolution; under a single global gradient-norm clip
that starves the remaining objectives, which is the same failure mode the
unnormalised $\mathcal{L}_{\text{smooth}}$ produced. Weighting each cell by the
map area it represents holds the term invariant.

\subsection{Toward Domain-Agnostic Multi-Region Observers}

The MRVP formulation and Director-CenterNet architecture are 
designed to be domain-agnostic. The structured game-state tensor can be 
replaced by any multi-channel spatial representation (e.g., MOBA minimap 
channels, drone swarm positional heatmaps, multi-CCTV occupancy grids). 
The four-loss framework remains applicable wherever observational data from
several human operators is available, since the ranked modes that supervise the
auxiliary regions are derived from disagreement among those operators rather
than from any game-specific quantity. Where only a single operator's data
exists, the modes collapse and the formulation degenerates to the single-region
case.

\subsection{Limitations}
\label{subsec:limitations}

\paragraph{Observer pool size.} With $U=5$ the support of a mode is an integer
in $\{1,\ldots,5\}$, so the ranking the method is built on has five levels and
the smallest auxiliary amplitude is $1/U = 0.2$. A mode carried by one
observer sits at the edge of what five samples separate from idiosyncratic
attention, and the confidence threshold has to be set below $0.2$ to admit it
at all. A larger pool would refine the ranking and make a support fraction,
rather than a count, the natural amplitude; it would also let the auxiliary
regions be supervised by modes that are individually better attested. We did
not collect one, and the results here should be read as what the ranking
recovers at the resolution the existing corpus provides.

\paragraph{Single-region accuracy.} At an unconstrained region budget the
Mask R-CNN + prior baseline remains ahead of the proposed method on IR
($0.684$ vs.\ $0.601$) and on OC3@0.5 ($0.592$ vs.\ $0.516$). The comparison at
a matched budget in Section~\ref{subsec:budget} reverses both, and we regard
that as the meaningful one, but it does not dispose of the first: a
practitioner who is content to display seven candidate regions and read only
the top one is better served by the baseline on these data. We have not closed
this gap and we do not know its cause. Head capacity is not it: widening every
head from $64$ to $256$ channels changes IR by $-0.017$, inside the
seed-to-seed band of $0.021$ we measure for this configuration. Heatmap
resolution remains the leading candidate and is under study. Until it
accounts for the gap, the contribution of this work should be read as being about
\emph{how many} regions a model emits and \emph{how they are ranked}, not about
the accuracy of the primary region.

\paragraph{Statistical power of the ablation.} The per-objective results in
Section~\ref{subsec:ablation} are single-seed, and a replicate of the full
configuration at a second seed differs by $0.021$ in IR, $0.015$ in OC3@0.5 and
$0.38$ in VD. Every difference between neighbouring configurations in both
ablation tables is smaller than the corresponding band, VD included. The
ablation therefore orders nothing, and no claim in this paper rests on it:
what the objectives are argued to do is argued from their construction and from
the aggregate comparisons, not from the lattice. Three seeds per cell are
required before those rows can be read at all, and we report them as measured
rather than omitting them so that the null result is on the record.

\paragraph{Held-out conditions.} All results are reported on the standard test
set, which fixes race matchup, map, and starting location. Earlier work
additionally evaluates transfer to unseen races, starting locations and maps
\cite{joo2023learning, anonymous2025kbrs}; we do not, and the accuracy reported
here should not be read as evidence about those conditions. The question is
also sharper for this formulation than for a single-region one, since the
auxiliary regions are supervised by minority modes whose spatial distribution
depends on map layout and on where each player starts. Whether the ranking
learned on one map transfers to another is therefore an open question and the
natural next experiment.

\paragraph{Track~B (zero human-GT setting).} Director-CenterNet currently
operates in Track~A (human observation data available for the main viewport).
Extending to Track~B — where no human viewport data exists and the system
must rely entirely on unsupervised saliency signals — is left as future work.

\paragraph{Temporal modelling.} $\mathcal{L}_{\text{smooth}}$ acts on 
pairs of consecutive frames. Incorporating longer-range temporal context via 
recurrent modules or temporal self-attention may further reduce thrashing 
for events that evolve over several seconds.

\paragraph{PIP count.} We fix $K=3$ in all experiments. Adapting $K$ 
dynamically beyond the confidence-threshold mechanism could improve coverage 
when more than two concurrent events occur simultaneously.

\paragraph{Live integration.} All experiments are conducted in offline 
replay mode. Real-time integration requires addressing the pipeline latency 
between game-state extraction, inference, and camera actuation.

%%====================================================================
\section{Conclusion}
\label{sec:conclusion}
%%====================================================================

We presented Director-CenterNet, a multi-region viewport prediction 
framework for automated Esports broadcasting. By reformulating automatic 
observing as Multi-Region Viewport Prediction and introducing four 
complementary training losses — Human Consensus Match, Ranked Mode Coverage,
Spatial Repulsion, and Trajectory Smoothness — Director-CenterNet predicts an
ordered set of presentation-agnostic regions from multi-channel game-state
tensors in a single forward pass. The Ranked Mode Coverage loss supervises the
auxiliary regions with the minority observer attention that previous work
discarded as noise, so no new form of annotation is required, while the
Spatial Repulsion and Trajectory Smoothness losses remove redundant overlap
and camera thrashing respectively.

Experiments on StarCraft: Brood War replays locate the contribution precisely.
A proposal-based detector emits several regions per frame already, so
multiplicity is not what distinguishes the formulations. What distinguishes
them is that the detector's region count moves in the wrong direction: it emits
fewer regions on the frames whose target has more modes, at every confidence
threshold we measured, because a per-proposal score is divided among boxes
exactly when the evidence is. No threshold reverses the sign, and raising one
does not improve the primary region — it only withholds answers, and withholds
them disproportionately on the multimodal frames the task is about. The heatmap
formulation reverses the sign, and it does so by virtue of how peaks are scored
and admitted rather than through the objectives we add on top of it. Held to
the same region budget, Director-CenterNet
leads on per-observer coverage and on count calibration; released from that
constraint, the strongest single-region baseline remains ahead on
primary-region overlap, and closing that remaining gap — most plausibly a
matter of heatmap resolution — is the immediate next step.

Future work will explore Track~B (zero human-GT) operation via unsupervised
saliency estimation, longer-range temporal context modelling, and extension
to MOBA games and non-game surveillance domains.

%%====================================================================
%% Declaration
%%====================================================================

\section*{CRediT authorship contribution statement}
\todo{Fill in author contributions following CRediT taxonomy.}

\section*{Declaration of competing interest}
The authors declare that they have no known competing financial interests 
or personal relationships that could have appeared to influence the work 
reported in this paper.

\section*{Acknowledgements}
\todo{Add funding acknowledgements.}

%%====================================================================
\appendix

\section{Mode Extraction Procedure}
\label{app:modes}

Algorithm-level detail for Section~\ref{subsec:modes}. Given the $U$ observer
viewports at frame $t$, we form the coverage map $a_t$ of
Eq.~\ref{eq:coverage}, smooth it with a Gaussian filter of width $\sigma_m$
under zero padding, and retain grid points that are local maxima of the
smoothed field under a $3{\times}3$ maximum filter and that exceed
$\theta \cdot \max \tilde{a}_t$. Candidates are visited in descending order of
$\tilde{a}_t$ and a candidate is kept only if it lies at least $D$ tiles from
every mode already kept. Each kept mode receives the support of
Eq.~\ref{eq:support}, and modes are ordered by support with ties broken by
smoothed response. At most $M$ modes are retained per frame.

The minimum separation $D$ is the parameter that decides what counts as two
distinct modes rather than one broad region of agreement, and it interacts
with the extraction stride: the non-maximum suppression applied to the
predicted heatmap must have a radius no larger than $D/2$ in tile units, or a
single mode yields several peaks within its own viewport and the region count
is inflated by duplicates.

\section{Hyperparameter Summary}
\label{app:hyperparams}

\begin{table}[h]
  \centering
  \caption{Hyperparameter settings used in all experiments unless 
  otherwise stated.}
  \label{tab:hyperparams}
  \begin{tabular}{lll}
    \toprule
    Hyperparameter & Value & Description \\
    \midrule
    Backbone         & ResNet-50  & Pretrained on ImageNet \\
    FPN stride $s$   & 4          & Finest FPN level; $32{\times}32$ heatmap \\
    Map size         & $128{\times}128$ tiles & Consumed at native resolution \\
    Viewport $(h,w)$ & $12{\times}20$ tiles & $3{\times}5$ heatmap cells \\
    Gaussian $\sigma$  & 4 tiles  & GT rendering, on the assigned cell (1 cell) \\
    Head width         & 64       & Hidden convs in each head \\
    Trainable stages   & all      & The heads read only the finest level \\
    Stacking $\Delta$  & 8 frames & Temporal stride \\
    Stacking $W$       & 3        & Number of past frames \\
    Mode $\sigma_m$    & 4 tiles  & Coverage-map smoothing \\
    Mode $\theta$      & 0.35     & Relative peak floor \\
    Mode $D$           & 12 tiles & Minimum mode separation \\
    Mode $M$           & 5        & Max modes retained per frame \\
    $\lambda_{\text{off}}$   & 1.0 & Offset loss weight \\
    $\lambda_{\text{sz}}$    & 0.1 & Size loss weight \\
    $\lambda_{\text{rmc}}$   & 0.5 & Ranked mode coverage weight \\
    $\lambda_{\text{rep}}$   & 0.3 & Repulsion loss weight \\
    $\lambda_{\text{sm}}$    & 0.1  & Smoothness loss weight \\
    $\Delta_s$       & 5 tiles    & Huber knee for $\mathcal{L}_{\text{smooth}}$ \\
    $E_0, E_1$       & 5, 10      & Smoothness warmup start / full \\
    $r$              & 2 cells    & Soft-centre window radius \\
    $\epsilon$       & 0.01       & Auxiliary support floor \\
    $K$              & 3          & Size of the predicted region set \\
    $\tau$           & 0.1        & Peak confidence threshold ($< 1/U$) \\
    Optimiser        & SGD        & momentum 0.9 \\
    Learning rate    & $5{\times}10^{-3}$ & StepLR (step 3, $\gamma$ 0.1) \\
    Weight decay     & $5{\times}10^{-4}$ & \\
    Batch size       & 16         & \\
    Epochs           & 30         & \\
    \bottomrule
  \end{tabular}
\end{table}

%%====================================================================
%% Bibliography
%%====================================================================

%% \todo{Add a references.bib entry for chen2018learning (Chen et al., learning
%% a director's style / reduced candidate set; cited in Sections 1 and 3 for
%% the first-versus-second-suggestion comparison that motivates BoK@delta).
%% The keys shi2016hierarchical, chen2011active, krause2008near, itti1998model,
%% wei2012geodesic, tatler2007central, feichtenhofer2017detect and
%% gandhi2014video are inherited from the previous draft and several are no
%% longer cited; prune them once the bibliography is regenerated.}

\bibliographystyle{elsarticle-num}
\bibliography{references}

\end{document}
